\documentclass[10pt]{article}

\usepackage{graphicx}
\usepackage[utf8]{inputenc}
\usepackage[spanish]{babel}
\usepackage{minted}
\usepackage{amsmath}
\usepackage[a4paper,top=2cm,bottom=2cm,left=3cm,right=3cm, marginparwidth=1.75cm]{geometry}
\usepackage{amsfonts}
\DeclareMathAlphabet{\pazocal}{OMS}{zplm}{m}{n}
\newcommand{\Lb}{\pazocal{L}}


\usemintedstyle{default}

\title{Práctica 3: APND, CBR, MT y REC}
\author{Miguel Ángel Dorado Maldonado}

\begin{document}
\maketitle

\section{Actividad 1: APND}

\[
    L = \{0^n1^{2n} : n > 0\}
\]

\subsection{Transiciones}

\begin{center}
    \centering
    \includegraphics[width=1\linewidth]{imageeee.png}
\end{center}

\subsection{Pruebas en JFLAP}

\begin{center}
    \centering
    \includegraphics[width=1\linewidth]{imageee.png}
\end{center}

\subsection{Definición en Python}

El APND se añade al archivo \texttt{pushdownautomata.json} con el nombre
\texttt{ej1\_0n1\_2n}:

\begin{minted}{json}
{
  "name": "ej1_0n1_2n",
  "hint": "APND para L={0^n1^2n : n>0}",
  "e.g.": [],
  "representation": {
    "K": ["q0", "q1"],
    "I": ["0", "1"],
    "S": ["A"],
    "s": "q0",
    "F": ["q1"],
    "t": [
      [["q0", "0", ""], ["q0", "AA"]],
      [["q0", "1", "A"], ["q1", ""]],
      [["q1", "1", "A"], ["q1", ""]]
    ]
  }
}
\end{minted}

Para comprobar el funcionamiento se usa el siguiente código:

\begin{minted}{python}
from python.automata import pushdown_automaton

pruebas = [
    ("011", "accept"),
    ("001111", "accept"),
    ("000111111", "accept"),
    ("01", "blocked"),
    ("00111", "blocked"),
    ("1111", "blocked"),
]

for cadena, esperado in pruebas:
    print("Cadena:", cadena)
    resultado = pushdown_automaton(
        "ej1_0n1_2n",
        cadena,
        desired_output=esperado
    )
    print("Resultado:", resultado)
    print("Esperado:", esperado)
    print()
\end{minted}

\section{Actividad 2: Máquina de Turing}

La entrada tiene la forma:

\[
    *|^{n+1}*|^{m+1}*
\]

La salida esperada es:

\[
    *|^{n+m+1}***
\]

\begin{center}
    \centering
    \includegraphics[width=1\linewidth]{sumaapdn.png}
\end{center}

\subsection{Definición en Python}

\begin{minted}{json}
{
  "name": "ej2_suma",
  "representation": {
    "matrix": [
      ["q0", "*", "l", "q1"],
      ["q0", "|", "|", "q0"],
      ["q1", "*", "|", "q2"],
      ["q1", "|", "l", "q1"],
      ["q2", "*", "l", "q3"],
      ["q2", "|", "r", "q2"],
      ["q3", "*", "l", "q4"],
      ["q3", "|", "*", "q3"],
      ["q4", "*", "h", "q4"],
      ["q4", "|", "*", "q4"]
    ]
  }
}
\end{minted}

\begin{minted}{python}
from python.turingmachine import turing_machine

pruebas = [
    "*|*|*",
    "*||*|*",
    "*|*||*",
    "*||*||*",
    "*|||*||*",
    "*|||*|||*",
]

for cinta in pruebas:
    print("Cinta:", cinta)
    cinta_final, posicion = turing_machine("ej2_suma", cinta, outputformat="none")
    print("Resultado:", "".join(cinta_final.content))
    print("Posicion:", posicion)
    print()
\end{minted}

\section{Actividad 3: Función recursiva}

La suma de tres valores se puede definir usando la suma de dos valores:

\[
    suma3(x,y,z) = suma(suma(x,y),z)
\]

Usando la notación de la librería:

\[
    addition(addition(\pi^3_1,\pi^3_2),\pi^3_3)
\]

\subsection{Evaluación a mano}

Tomamos como ejemplo:

\[
    suma3(2,3,4)
\]

\[
\begin{aligned}
suma3(2,3,4) &= suma(suma(2,3),4) \\
              &= suma(5,4) \\
              &= 9
\end{aligned}
\]

\subsection{Ejecución en Python}

\begin{minted}{python}
from python.recursivefunctions import eval_rec_function

suma3 = "addition(addition(π^3_1,π^3_2),π^3_3)"

print(eval_rec_function(suma3, 2, 3, 4))
\end{minted}

El resultado obtenido es:

\begin{minted}{text}
9
\end{minted}


\end{document}
